% ID: SEL-D0024-OA-E01
% ESTADO: validada. S025, cotejada con examen y criterios.
Sean $x$ e $y$ los envases de los tipos primero y segundo.
Pasamos las disponibilidades a gramos y dividimos cada restricción
entre 10. El modelo queda
\[
3x+4y\le5200,\quad 4x+3y\le4600,\quad 2x+y\le2000,
\qquad x,y\ge0.
\]
Maximizamos $G=1.5x+y$, en euros. En los ejes los límites más
restrictivos son $x\le1000$ e $y\le1300$. Las dos primeras
fronteras se cortan en $(400,1000)$, y la segunda y tercera
en $(700,600)$. La intersección de la primera y la tercera
es $(560,880)$, que viola la segunda restricción.
\[
\begin{array}{c|rrrrr}
(x,y)&(0,0)&(1000,0)&(700,600)&(400,1000)&(0,1300)\\\hline
G(x,y)&0&1500&1650&1600&1300
\end{array}
\]

\begin{center}
\setlength{\unitlength}{3.5mm}
\begin{picture}(20,17)(-4,-2)
\qbezier(0,0)(6,0)(12,0)\put(12.2,0){$x$}
\qbezier(0,0)(0,7)(0,14)\put(0,14.2){$y$}
\qbezier(0,13)(2,11.5)(4,10)
\qbezier(4,10)(5.5,8)(7,6)
\qbezier(7,6)(8.5,3)(10,0)
\put(-0.3,13){\makebox(0,0)[r]{$(0,1300)$}}
\put(4.4,10.4){$(400,1000)$}
\put(7.4,6.3){$(700,600)$}
\put(10,-1){\makebox(0,0)[t]{$(1000,0)$}}
\end{picture}
\end{center}

La función lineal alcanza su máximo en un vértice. El mayor valor
es 1650, y el vértice correspondiente tiene coordenadas enteras,
por lo que también es óptimo entre cantidades enteras.
Se deben fabricar \textbf{700 envases del primer tipo y 600 del segundo}.
El importe de la venta es $\boxed{1650}$ euros. El enunciado llama
ganancia a ese importe y sólo facilita precios de venta.
